\section{E-commerce Sales Business}
\label{sec:2-1}

The business sells products through an e-commerce transaction system, and sales generate revenue from customer purchases. At a high level:
\begin{equation*}
    \text{Profit} = \text{Sales Revenue} - \text{Business Costs}
\end{equation*}
Costs include product and operating costs, marketing and contact costs, and the discounts given as retention incentives.

\section{Revenue from New and Existing Customers}
\label{sec:2-2}

Sales revenue comes from new customers making their first purchase and from existing customers making repeat purchases. In this dataset almost all revenue comes from customers who buy more than once (Section~\ref{sec:4-4}), so keeping existing customers buying is where the value is.

\section{Customer Churn as a Business Problem}
\label{sec:2-3}

When existing customers stop purchasing, the business loses their future sales. Retail has no contracts, so customers never cancel: churn is inferred from a period without purchases (Section~\ref{sec:4-5}) and must be anticipated while the customer is still active.

\section{Budget and Capacity Constraints}
\label{sec:2-4}

Retention activities require resources. The retailer works with a \textbf{fixed weekly discount budget} and a limited capacity to contact customers, so it cannot send the same offer to every active customer. Retention resources must be allocated selectively.

\section{Business Objectives}
\label{sec:2-5}

\textbf{The decision.} Each Monday, the retailer must choose whom to contact within a fixed discount budget. Every offer has one of two outcomes:

\begin{itemize}
    \item \textbf{A purchase we might recover.} A customer who would otherwise stop buying may return because of the offer.
    \item \textbf{A purchase we might subsidise.} A customer who would buy anyway may use the same discount.
\end{itemize}

\textbf{Decision criterion:} maximise expected incremental profit relative to sending no offer. The specific objectives are:

\begin{itemize}
    \item \textbf{Identify at-risk customers:} estimate, for every active customer, the probability of making no purchase within a predefined period.
    \item \textbf{Understand purchasing behavior:} investigate which characteristics -- purchase frequency, monetary value, time since the last purchase -- are associated with churn risk.
    \item \textbf{Choose the offer list:} rank customers so that a fixed weekly discount budget goes where it adds the most margin, and state how to verify this in a pilot.
\end{itemize}

\section{Success Criteria}
\label{sec:2-6}

Success is defined \textbf{relative to simple alternatives} rather than by arbitrary absolute thresholds:

\begin{center}
\begin{tabular}{@{}p{3cm}p{11.5cm}@{}}
    \toprule
    Level & Criterion (evaluated on the unseen test period) \\
    \midrule
    Predictive & Final model AUC exceeds random targeting and the recency-rule baseline. \\
    Targeting & At the same expected weekly discount spend, the list ranked by the model reaches more customers who would otherwise stop buying than the recency rule. \\
    Stability & Weekly AUC does not collapse on individual test weeks, and the weekly score distribution stays within a moderate shift of the validation reference (CSI $< 0.25$). \\
    Business scenario & Under the base scenario, the simulated net margin of the model-ranked list is positive and higher than for the recency rule at the same expected discount spend, in most test weeks. \\
    \bottomrule
\end{tabular}
\end{center}

\section{From Business Objective to Analytical Problem}
\label{sec:2-7}

Increasing profit from existing customers requires recovering purchases that would otherwise be lost while giving away as little margin as possible to customers who would buy anyway. The analytical direction is therefore to estimate each active customer's churn probability, translate it into the expected cost and margin of an offer, and choose the weekly offer list within the discount budget.

\section{Customer Retention Strategy}
\label{sec:2-8}

\begin{center}
\begin{tabular}{@{}>{\centering\arraybackslash}p{14cm}@{}}
    \textbf{Eligible Existing Customers} -- active on the scoring Monday \\
    $\downarrow$ \\
    \textbf{Churn Prediction} -- estimate the probability of no purchase in the next 90 days \\
    $\downarrow$ \\
    \textbf{Risk Ranking} -- rank customers by predicted churn risk \\
    $\downarrow$ \\
    \textbf{Weekly Discount Budget} -- offer down the ranking until the expected discount cost reaches the budget \\
    $\downarrow$ \\
    \textbf{Voucher List and Pilot} -- send the offer, keep a no-offer control group, measure the result \\
\end{tabular}
\end{center}

\section{Implementation Roadmap}
\label{sec:2-9}

The approach is implemented as a sequence of steps, each covered by one chapter and one or two notebooks in \texttt{notebooks/}:

\begin{itemize}
    \item \textbf{Understand the data} (Chapter~\ref{ch:data-understanding}; \texttt{01\_data\_understanding}, \texttt{02\_customer\_segments}) -- revenue and customers, purchase intervals and customer segments, which inform the prediction horizon $H$ and the activity threshold $R_{\max}$.
    \item \textbf{Prepare customer snapshots} (Chapter~\ref{ch:data-preparation}; \texttt{03\_data\_preparation}) -- turn transactions into Customer $\times$ Snapshot Date rows with point-in-time features, churn labels and an embargoed temporal split.
    \item \textbf{Model churn risk} (Chapter~\ref{ch:modeling}; \texttt{04\_modeling}) -- compare baselines, training rows, feature sets and algorithms, tune the final model and evaluate it once on the test weeks.
    \item \textbf{Evaluate the business impact} (Chapter~\ref{ch:evaluation}; \texttt{05\_scoring}, \texttt{06\_business\_impact}) -- score every customer-week and compare ways of choosing customers at the same discount budget.
    \item \textbf{Decide and monitor} (Chapter~\ref{ch:targeting-policy}; \texttt{07\_action\_table}, \texttt{dashboard/}) -- the weekly voucher list, the dashboard and the monitoring of the score, as a proposed operational design.
\end{itemize}
